\subsection{忠实平坦环上的线性方程组}

设 B 是环，A 是 B 的子环。若有序偶 (A, B) 满足下列条件，则称它具有线性扩张性质：

\begin{enumerate}
\def\labelenumi{(\Alph{enumi})}
\setcounter{enumi}{4}
\tightlist
\item
  线性方程组的每个由 B 的元素组成的解 \((y_{k})_{1\leqslant k\leqslant n}\) ，
\end{enumerate}

\[
\sum_ {k = 1} ^ {n} y _ {k} c _ {k \mathrm{l}} = d _ {i} \quad (1 \leqslant i \leqslant m)\tag{3}
\]

其系数 \(c_{kl}\) 与右端 \(d_{i}\) 都属于 A，它形如

\[
y _ {k} = x _ {k} + \sum_ {j = 1} ^ {p} b _ {j} z _ {j k} \quad (1 \leqslant k \leqslant n)\tag{4}
\]

其中 \((x_{k})\) 是 (3) 的由 A 的元素组成的解，诸 \(b_{j}\) 属于 B，且每个 \((z_{jk})_{1\leq k\leq n}\) 是与 (3) 相伴的齐次线性方程组的由 A 的元素组成的解。

命题 13. 设 A 是环 B 的子环。为使有序偶 (A, B) 具有线性扩张性质，必须且只需 B 是忠实平坦 A-模。

条件是充分的。因为，由于 B 是平坦 A-模，与 (3) 相伴的齐次线性方程组的每个在 B 中取值的解，都是以 B 为系数、对由 A 中元素组成的解所作的线性组合（ \(\S\) 2 第 11 目，命题 13 的推论 2）。于是问题归结为证明：若 (3) 在 B 中有解，则它在 A 中也有一个解。现在若置

\[
c _ {k} = \left(c _ {k l}\right) _ {1 \leqslant i \leqslant m} \in \mathrm{A} _ {s} ^ {m}, \quad d = (d _ {i}) \in \mathrm{A} _ {s} ^ {m},
\]

方程组 (3) 等价于方程 \(\sum_{k=1}^{n} y_k \otimes c_k = 1 \otimes d\) ，后者在 \(\mathbf{B} \otimes_{\mathbf{A}} \mathbf{A}_s^m = \mathbf{B}_s^m\) 中成立。换言之，若 \(\mathbf{M}\) 是 \(\mathbf{A}_s^m\) 的由诸 \(c_k\) ( \(1 \leqslant k \leqslant n\) ) 生成的子 A-模，则 (3) 的解 \((y_k)\) 的存在性等价于（取第 5 目中所作的等同）关系 \(d \in \mathbf{BM} \cap \mathbf{A}_s^m\) ；但由于 \(\mathbf{BM} \cap \mathbf{A}_s^m = \mathbf{M}\) （第 5 目，命题 10，(ii)），它蕴含 \(d \in \mathbf{M}\) ，即存在解 \((x_k)\) ，它由 \(\mathbf{A}\) 中元素组成，是方程组 (3) 的解。

条件是必要的。因为假设 (A, B) 具有线性扩张性质；我们已知 B 是平坦右 A-模（ \(\S\) 2 第 11 目，命题 13 的推论 2）；我们证明，对 A 的每个左理想 a， \(Ba \cap A = a\) ，这表明 B 是忠实平坦右 A-模（第 5 目，命题 9，(d)）。现在设 \(x \in Ba \cap A\) ；由假设存在 \(y_i \in B\) 与 \(a_i \in a\) 使 \(\sum_{i} y_i a_i = x\) ；把性质 (E) 应用于这个系数与右端都在 A 中的线性方程，可知存在 \(x_i \in A\) 使 \(x = \sum_{i} x_i a_i\) ，因此 \(x \in a\) 。

